\section{Load/Store Unit and Memory Access}

\subsection{Why RISC-V is load/store}

RISC-V is a load/store architecture. Arithmetic instructions operate on registers, not directly on memory. To use a memory value, software must first load it into a register. To update memory, software must store a register value.

\begin{lstlisting}[language={}]
lw   x5, 0(x10)   # x5 = memory[x10 + 0]
add  x6, x5, x7   # arithmetic uses registers
sw   x6, 4(x10)   # memory[x10 + 4] = x6
\end{lstlisting}

\subsection{Effective address calculation}

Loads and stores compute an effective address:

\begin{lstlisting}[language={}]
effective_address = rs1 + sign_extended_immediate
\end{lstlisting}

The ALU performs this addition. In simple cores, this address goes directly to \texttt{dmem.v}. In SoC-oriented cores, it goes to an external data bus. In MMU variants, it may first be translated from virtual to physical.

\subsection{Load sizes and sign extension}

RISC-V loads can read different sizes:

\begin{longtable}{|>{\RaggedRight\arraybackslash}p{0.28\textwidth}|>{\RaggedRight\arraybackslash}p{0.32\textwidth}|>{\RaggedRight\arraybackslash}p{0.30\textwidth}|}
\hline
\textbf{Instruction} & \textbf{Size} & \textbf{Extension} \\
\hline
\endfirsthead
\hline
\textbf{Instruction} & \textbf{Size} & \textbf{Extension} \\
\hline
\endhead
\texttt{lb} & 8-bit byte & sign-extend \\
\hline
\texttt{lbu} & 8-bit byte & zero-extend \\
\hline
\texttt{lh} & 16-bit halfword & sign-extend \\
\hline
\texttt{lhu} & 16-bit halfword & zero-extend \\
\hline
\texttt{lw} & 32-bit word & no extension needed \\
\hline
\end{longtable}

The memory block or CPU memory stage must select the correct byte/halfword lane and extend it to 32 bits.

\subsection{Store sizes and byte masks}

Stores can write bytes, halfwords, or words:

\begin{longtable}{|>{\RaggedRight\arraybackslash}p{0.38\textwidth}|>{\RaggedRight\arraybackslash}p{0.52\textwidth}|}
\hline
\textbf{Instruction} & \textbf{Size} \\
\hline
\endfirsthead
\hline
\textbf{Instruction} & \textbf{Size} \\
\hline
\endhead
\texttt{sb} & 8-bit byte \\
\hline
\texttt{sh} & 16-bit halfword \\
\hline
\texttt{sw} & 32-bit word \\
\hline
\end{longtable}

The multi-cycle memory design uses byte write masks such as \texttt{4'b0001 << byte\_offset} for byte stores and similar masks for halfword/word stores. This is important for block RAM and bus integration.

\subsection{Memory-mapped I/O}

In \texttt{cpu\_core.v}, the CPU does not own data memory directly. Instead it exposes:

\begin{lstlisting}[language={}]
dmem_addr
dmem_wdata
dmem_we
dmem_funct3
dmem_rdata
\end{lstlisting}

A SoC wrapper decodes the address and decides whether the access targets RAM, UART, timer, syscon, or another peripheral. This is memory-mapped I/O: software uses normal load/store instructions, but some addresses control hardware devices instead of RAM.

\bigskip
